\section{Graded rings}
\label{section-graded}

\noindent
A {\it graded ring} will be for us a ring $S$ endowed
with a direct sum decomposition $S = \bigoplus_{d \geq 0} S_d$
of the underlying abelian group such that $S_d \cdot S_e \subset S_{d + e}$.
Note that we do not allow nonzero elements in negative degrees.
The {\it irrelevant ideal} is the ideal $S_{+} = \bigoplus_{d > 0} S_d$.
A {\it graded module}
will be an $S$-module $M$ endowed with a direct sum decomposition
$M = \bigoplus_{n\in \mathbf{Z}} M_n$ of the underlying abelian group
such that $S_d \cdot M_e \subset M_{d + e}$. Note that for modules we do allow
nonzero elements in negative degrees.
We think of $S$ as a graded $S$-module by setting $S_{-k} = (0)$
for $k > 0$. An element $x$ (resp.\ $f$) of $M$ (resp.\ $S$) is called
{\it homogeneous}
if $x \in M_d$ (resp.\ $f \in S_d$) for some $d$.
A {\it map of graded $S$-modules} is a map of $S$-modules
$\varphi : M \to M'$ such that $\varphi(M_d) \subset M'_d$.
We do not allow maps to shift degrees. Let us denote
$\text{GrHom}_0(M, N)$ the $S_0$-module of homomorphisms
of graded modules from $M$ to $N$.

\medskip\noindent
At this point there are the notions of graded ideal,
graded quotient ring, graded submodule, graded quotient module,
graded tensor product, etc. We leave it to the reader to find the
relevant definitions, and lemmas. For example: A short exact sequence
of graded modules is short exact in every degree.

\medskip\noindent
Given a graded ring $S$, a graded $S$-module $M$ and $n \in \mathbf{Z}$
we denote $M(n)$ the graded $S$-module with $M(n)_d = M_{n + d}$.
This is called the {\it twist of $M$ by $n$}. In particular we get
modules $S(n)$, $n \in \mathbf{Z}$ which will play an important
role in the study of projective schemes. There are some obvious
functorial isomorphisms such as
$(M \oplus N)(n) = M(n) \oplus N(n)$,
$(M \otimes_S N)(n) = M \otimes_S N(n) = M(n) \otimes_S N$.
In addition we can define a graded $S$-module structure on
the $S_0$-module
$$
\text{GrHom}(M, N) =
\bigoplus\nolimits_{n \in \mathbf{Z}} \text{GrHom}_n(M, N),
\quad
\text{GrHom}_n(M, N) = \text{GrHom}_0(M, N(n)).
$$
We omit the definition of the multiplication.

\begin{lemma}
\label{lemma-graded-NAK}
Let $S$ be a graded ring. Let $M$ be a graded $S$-module.
\begin{enumerate}
\item If $S_+M = M$ and $M$ is finite, then $M = 0$.
\item If $N, N' \subset M$ are graded submodules,
$M = N + S_+N'$, and $N'$ is finite, then $M = N$.
\item If $N \to M$ is a map of graded modules, $N/S_+N \to M/S_+M$
is surjective, and $M$ is finite, then $N \to M$ is surjective.
\item If $x_1, \ldots, x_n \in M$ are homogeneous and generate $M/S_+M$
and $M$ is finite, then $x_1, \ldots, x_n$ generate $M$.
\end{enumerate}
\end{lemma}

\begin{proof}
Proof of (1). Choose generators $y_1, \ldots, y_r$ of $M$ over $S$.
We may assume that $y_i$ is homogeneous of degree $d_i$. After
renumbering we may assume $d_r = \min(d_i)$. Then the condition that
$S_+M = M$ implies $y_r = 0$. Hence $M = 0$ by induction on $r$.
Part (2) follows by applying (1) to $M/N$. Part (3) follows by
applying (2) to the submodules $\Im(N \to M)$ and $M$.
Part (4) follows by applying (3) to the module map
$\bigoplus S(-d_i) \to M$, $(s_1, \ldots, s_n) \mapsto \sum s_i x_i$.
\end{proof}

\noindent
Let $S$ be a graded ring. Let $d \geq 1$ be an integer.
We set $S^{(d)} = \bigoplus_{n \geq 0} S_{nd}$. We think of
$S^{(d)}$ as a graded ring with degree $n$ summand
$(S^{(d)})_n = S_{nd}$. Given a graded $S$-module $M$ we
can similarly consider $M^{(d)} = \bigoplus_{n \in \mathbf{Z}} M_{nd}$
which is a graded $S^{(d)}$-module.

\begin{lemma}
\label{lemma-uple-generated-degree-1}
Let $S$ be a graded ring, which is finitely generated over $S_0$.
Then for all sufficiently divisible $d$ the algebra
$S^{(d)}$ is generated in degree $1$ over $S_0$.
\end{lemma}

\begin{proof}
Say $S$ is generated by $f_1, \ldots, f_r \in S$ over $S_0$.
After replacing $f_i$ by their homogeneous parts, we may assume
$f_i$ is homogeneous of degree $d_i > 0$. Then any element of
$S_n$ is a linear combination with coefficients in $S_0$ of monomials
$f_1^{e_1} \ldots f_r^{e_r}$ with $\sum e_i d_i = n$.
Let $m$ be a multiple of $\text{lcm}(d_i)$. For any $N \geq r$ if
$$
\sum e_i d_i = N m
$$
then for some $i$ we have $e_i \geq m/d_i$ by an elementary argument.
Hence every monomial of degree $N m$ is a product of a monomial
of degree $m$, namely $f_i^{m/d_i}$, and a monomial of degree $(N - 1)m$.
It follows that any monomial of degree $nrm$ with $n \geq 2$
is a product of monomials of degree $rm$. Thus $S^{(rm)}$ is generated
in degree $1$ over $S_0$.
\end{proof}

\begin{lemma}
\label{lemma-integral-closure-graded}
\begin{reference}
\cite[Theorem 2.3.2]{Huneke-Swanson}
\end{reference}
Let $R \to S$ be a homomorphism of graded rings.
Let $S' \subset S$ be the integral closure of $R$ in $S$.
Then
$$
S' = \bigoplus\nolimits_{d \geq 0} S' \cap S_d,
$$
i.e., $S'$ is a graded $R$-subalgebra of $S$.
\end{lemma}

\begin{proof}
We have to show the following: If
$s = s_n + s_{n + 1} + \ldots + s_m \in S'$, then each homogeneous
part $s_j \in S'$. We will prove this by induction on $m - n$ over
all homomorphisms $R \to S$ of graded rings. First note that it
is immediate that $s_0$ is integral over $R_0$ (hence over $R$) as
there is a ring map $S \to S_0$ compatible with the ring map $R \to R_0$.
Thus, after replacing $s$ by $s - s_0$, we may assume $n > 0$. Consider the
extension of graded rings $R[t, t^{-1}] \to S[t, t^{-1}]$ where
$t$ has degree $0$. There is a commutative diagram
$$
\xymatrix{
S[t, t^{-1}] \ar[rr]_{s \mapsto t^{\deg(s)}s} & & S[t, t^{-1}] \\
R[t, t^{-1}] \ar[u] \ar[rr]^{r \mapsto t^{\deg(r)}r} & &  R[t, t^{-1}] \ar[u]
}
$$
where the horizontal maps are ring automorphisms. Hence the integral
closure $C$ of $S[t, t^{-1}]$ over $R[t, t^{-1}]$ maps into itself.
Thus we see that
$$
t^m(s_n + s_{n + 1} + \ldots + s_m) -
(t^ns_n + t^{n + 1}s_{n + 1} + \ldots + t^ms_m) \in C
$$
which implies by induction hypothesis that each $(t^m - t^i)s_i \in C$
for $i = n, \ldots, m - 1$. Note that for any ring $A$ and $m > i \geq n > 0$
we have $A[t, t^{-1}]/(t^m - t^i - 1) \cong A[t]/(t^m - t^i - 1) \supset A$
because $t(t^{m - 1} - t^{i - 1}) = 1$ in $A[t]/(t^m - t^i - 1)$.
Since $t^m - t^i$ maps to $1$ we see the image of $s_i$ in the ring
$S[t]/(t^m - t^i - 1)$ is integral over $R[t]/(t^m - t^i - 1)$ for
$i = n, \ldots, m - 1$. Since $R \to R[t]/(t^m - t^i - 1)$ is finite
we see that $s_i$ is integral over $R$ by transitivity, see
Lemma \ref{lemma-integral-transitive}.
Finally, we also conclude that $s_m = s - \sum_{i = n, \ldots, m - 1} s_i$
is integral over $R$.
\end{proof}








